\section*{Analysis of Features 5 and 6}

\subsection*{What do you observe?}
Observing the scatter plot (\texttt{features\_5\_6\_scatter.png}) and the histograms (\texttt{feature\_5\_histogram.png} and \texttt{feature\_6\_histogram.png}), we can analyze the distribution of the "Fake" (label 0) and "Genuine" (label 1) classes. [Insert visual observation here: e.g., Discuss the general spread of the data, whether one class has a much wider range than the other, and any overarching patterns].

\subsection*{Do the classes overlap? If so, where?}
Yes, there is overlap between the classes. According to the computed data in \texttt{lab02\_project\_feature\_overlap\_ranges.csv}:
\begin{itemize}
    \item \textbf{Feature 5} overlaps between [Insert Min Overlap] and [Insert Max Overlap] (an overlap width of [Insert Overlap Width]).
    \item \textbf{Feature 6} overlaps between [Insert Min Overlap] and [Insert Max Overlap] (an overlap width of [Insert Overlap Width]).
\end{itemize}

\subsection*{How many modes are evident from the histograms?}
Based on the visual inspection of the density histograms:
\begin{itemize}
    \item \textbf{Feature 5:} The "Fake" distribution exhibits [Insert Number] mode(s) (peaks), and the "Genuine" distribution exhibits [Insert Number] mode(s).
    \item \textbf{Feature 6:} The "Fake" distribution exhibits [Insert Number] mode(s), and the "Genuine" distribution exhibits [Insert Number] mode(s).
\end{itemize}

\subsection*{How many clusters can you notice from the scatter plots for each class?}
Examining the scatter plot of Feature 5 vs Feature 6 (\texttt{features\_5\_6\_scatter.png}):
\begin{itemize}
    \item \textbf{Fake Class (Label 0):} There appear to be [Insert Number, e.g., one distinct cluster / multiple separated clusters].
    \item \textbf{Genuine Class (Label 1):} There appear to be [Insert Number] discernible cluster(s). [Insert any additional notes on cluster density or shape].
\end{itemize}
